\documentclass[10pt,a4paper]{amsart}
\usepackage[margin=19mm]{geometry}
\usepackage{amsmath,amssymb,mathtools}
\usepackage[scr=rsfs,cal=euler]{mathalfa}
\usepackage{xfrac}
\usepackage[unicode,hidelinks,bookmarks=false]{hyperref}
\newcommand{\ie}{i.e.\ }
\newcommand{\cf}{cf.\ }
\newcommand{\resp}{respectively\ }
\newcommand{\Subsection}{\subsection}
\providecommand{\nobreakdash}{\nobreak\hbox{-}}
\setlength{\emergencystretch}{2em}
\multlinegap0pt
\usepackage{amssymb}
\usepackage{mathtools}
\usepackage{bm}
\usepackage{xcolor}
\usepackage{tikz}
\newtheorem{proposition}{Proposition}
\newcommand{\C}{\mathbb{C}}
\newcommand{\Her}[2][\C]{\operatorname{Her}_{#2}(#1)}
\newcommand{\Z}{\mathbb{Z}}
\title{Graph Quantum Magic Squares and Free Spectrahedra}
\author{Francesca La Piana}
\date{Source: arXiv:2512.08797v3 (CC BY 4.0)}
\begin{document}
\maketitle

\noindent\emph{Source and licence.} Graph Quantum Magic Squares and Free Spectrahedra. Version arXiv:2512.08797v3 (CC BY 4.0), distributed by arXiv under the Creative Commons Attribution 4.0 International licence, http://creativecommons.org/licenses/by/4.0/, as recorded in the arXiv licence metadata for this exact version and shown on the abstract page https://arxiv.org/abs/2512.08797v3. Full licence notice: \texttt{licenses/arxiv-2512.08797-CC-BY-4.0.txt}.\par\smallskip

\noindent\textit{Editorial scope.} Counted unit: the article's proposition prop:Z4-averaging (source0.tex lines 928--965). The article's complete original proof of it is delivered with it (source0.tex lines 967--1070, one \texttt{proof} environment of 104 lines). No mathematics is written, repaired or re-derived: the statement and the proof are the article's own lines, in the article's own notation, and no retained line is substituted or re-flowed.  No further source-local context is needed: every cross-reference of every retained line resolves inside the two delivered spans.  Nothing was omitted from any retained span, so the package carries no omission record.  No retained line contains a citation, so the wrapper carries no bibliography and no bibliography entry.  No retained line contains a figure environment, an image-inclusion command or drawing code, so no image is delivered and the package declares no asset dependency.  Editorial formatting only, in addition: an AMS \texttt{amsart} preamble that restates the article's own declarations and macros used by the retained lines, namely the environments \texttt{proposition} on separate counters, together with the standard packages those lines need.  The retained environments print in this document as Proposition 1 in order of appearance, whereas the article prints the same items with the numbering of its own full document.  The complete native source of the article as distributed in the arXiv e-print for this exact version is bundled unchanged as \texttt{sources/190691/source0.tex}.
\begin{proposition}\label{prop:Z4-averaging}
Let \(A=\bigl(A_{ij}\bigr)_{i,j=1}^4 \in \operatorname{Mat}_4(\Her{s})\) be a \(4\times 4\) quantum magic square of internal size \(s\).
Let \(P_{C_4}\) be the \(4\times 4\) permutation matrix
\[
P_{C_4} e_i = e_{\,i+1 \pmod 4}, \qquad i=1,\dots,4,
\]
that is,
\[
P_{C_4} =
\begin{pmatrix}
0 & 0 & 0 & 1\\
1 & 0 & 0 & 0\\
0 & 1 & 0 & 0\\
0 & 0 & 1 & 0
\end{pmatrix}.
\]
Define
\[
\Phi(A) := \frac{1}{4}\sum_{k=0}^{3} 
\bigl(P_{C_4}^{k}\otimes I_s\bigr)\,
A\,
\bigl(P_{C_4}^{k}\otimes I_s\bigr)^{*}.
\]
Then:
\begin{enumerate}
    \item[(i)] \(\Phi(A)\) is a quantum magic square;
    \item[(ii)] \(\Phi(A)\) commutes with the adjacency matrix of \(C_4\), i.e.
    \[
    (A_{C_4}\otimes I_s)\, \Phi(A) = \Phi(A)\, (A_{C_4}\otimes I_s);
    \]
    \item[(iii)]  we have (entrywise):
    \[
    \Phi(A)_{ij}
      = \frac{1}{4}\sum_{k=0}^3 A_{\,i+k,\; j+k},
    \qquad (\text{indices } i,j \text{ mod } 4).
    \]
\end{enumerate}
\end{proposition}

\begin{proof}
Let us set \(\widetilde P_{C_4} := P_{C_4}\otimes I_s\).  
Then \(\widetilde P_{C_4}\) is unitary, so \(\widetilde P^*_{C_4}=\widetilde P^{-1}_{C_4}\) and,
for every \(k \in \Z\),
\[
\bigl(\widetilde P_{C_4}^{\,k}\bigr)^*
= \bigl(\widetilde P_{C_4}^*\bigr)^{k}
= \widetilde P_{C_4}^{-k}.
\]
Moreover \((P_{C_4})^4 = I_4\) implies \(\widetilde P_{C_4}^{\,4} = I_4\otimes I_s\).


\smallskip

(ii) 
We start showing that \(\Phi(A)\) commutes with \(\widetilde P_{C_4}\).
By definition of \(\Phi\) and linearity we have
\begin{align*}
\widetilde P_{C_4}\,\Phi(A)\,\widetilde P_{C_4}^{-1}
&= \widetilde P_{C_4} \left( \frac{1}{4}\sum_{k=0}^{3} 
\widetilde P_{C_4}^{\,k} A \bigl(\widetilde P_{C_4}^{\,k}\bigr)^{*} \right) \widetilde P_{C_4}^{-1} \\
&= \frac14 \sum_{k=0}^3 
\widetilde P_{C_4}^{\,k+1} A \widetilde P_{C_4}^{-(k+1)}.
\end{align*}
Changing the index \(m=k+1 \pmod 4\), the sum becomes
\[
\frac14 \sum_{m=0}^3 \widetilde P_{C_4}^{\,m} A \widetilde P_{C_4}^{-m} = \Phi(A),
\]
so \(\widetilde P_{C_4}\,\Phi(A)\,\widetilde P_{C_4}^{-1}=\Phi(A)\), i.e. 
\begin{equation}\label{comm_P_tilde}
    \widetilde P_{C_4}\,\Phi(A)=\Phi(A)\,\widetilde P_{C_4}.
\end{equation}

The adjacency matrix of the undirected cycle is
\[
A_{C_4} = P_{C_4} + P_{C_4}^*,
\qquad
A_{C_4}\otimes I_s
= \widetilde P_{C_4} + \widetilde P_{C_4}^{*}.
\]
Since \(\Phi(A)\) is a convex combination of conjugates of \(A\), if \(A\) is Hermitian \(\Phi(A)\) is too. Taking adjoints in \eqref{comm_P_tilde}, we get
\[
\Phi(A)\,\widetilde P_{C_4}^{*} = \widetilde P_{C_4}^{*} \,\Phi(A).
\]
So \(\Phi(A)\) commutes with both \(\widetilde P_{C_4}\) and \(\widetilde P_{C_4}^{*}\), and so does their sum:
\[
(A_{C_4}\otimes I_s)\,\Phi(A)
= \Phi(A)\,(A_{C_4}\otimes I_s).
\]
This proves (ii).

\smallskip

(iii)
For a block matrix \(M\), the \((i,j)\)-th block is
\(M_{ij}=(e_i^*\otimes I_s)\,M\,(e_j\otimes I_s)\).

For a fixed \(k\), we have
\begin{align*}
\bigl(\widetilde P_{C_4}^{\,k} A \widetilde P_{C_4}^{-k}\bigr)_{ij}
&= (e_i^*\otimes I_s)\, \widetilde P_{C_4}^{\,k} A \widetilde P_{C_4}^{-k}\, (e_j\otimes I_s) \\
&= \bigl((e_i^* P_{C_4}^{\,k})\otimes I_s\bigr)\, A \,
   \bigl((P_{C_4}^{-k} e_j)\otimes I_s\bigr).
\end{align*}
Using the equalities \(P_{C_4}^{-k}e_j=e_{\,j-k}\) and
\(e_i^* P_{C_4}^{\,k} = (P_{C_4}^{-k}e_i)^* = e_{\,i-k}^*\), we get
\[
\bigl(\widetilde P_{C_4}^{\,k} A \widetilde P_{C_4}^{-k}\bigr)_{ij}
= (e_{\,i-k}^* \otimes I_s)\, A\, (e_{\,j-k}\otimes I_s)
= A_{\,i-k,\; j-k}.
\]
Therefore,
\[
\Phi(A)_{ij}
= \left(\frac14 \sum_{k=0}^3 \widetilde P_{C_4}^{\,k} A \widetilde P_{C_4}^{-k}\right)_{ij}
= \frac14 \sum_{k=0}^3 A_{\,i-k,\; j-k}.
\]
Since the indices are taken modulo \(4\), summing over \(k\) is equivalent to sum over \(-k\), and then we can write
\[
\Phi(A)_{ij}
= \frac14 \sum_{k=0}^3 A_{\,i+k,\; j+k}
\]
proving (iii).

\smallskip

(i)
Recall that a quantum magic square \(A\) satisfies:
\begin{enumerate}
\item[(a)] \(A_{ij} \succeq 0\) for all \(i,j\);
\item[(b)] \(\sum_{j=1}^4 A_{ij} = I_s\) for each row \(i\);
\item[(c)] \(\sum_{i=1}^4 A_{ij} = I_s\) for each column \(j\).
\end{enumerate}
Conjugation by \(\widetilde P_{C_4}^k\) just permutes rows and columns (and hence the
blocks) of \(A\), so each \(\widetilde P_{C_4}^{\,k} A \widetilde P_{C_4}^{-k}\) is again a QMS:
the blocks remain positive semidefinite and the row/column sums stay equal to \(I_s\).

\(\Phi(A)\) is a convex combination of these matrices, so:
\begin{itemize}
\item each block \(\Phi(A)_{ij}\) is a convex combination of positive semidefinite blocks, hence \(\Phi(A)_{ij}\succeq 0\);
\item the row-sums and column-sums of \(\Phi(A)\) are averages of \(I_s\), hence still equal to \(I_s\).
\end{itemize}
Therefore, since \(\Phi(A)\) satisfies (a), (b), (c), it is a quantum magic square and this proves (i).
\end{proof}

\end{document}
